\section{OPTICAL PLATFORMS AND ARCHITECTURE CHOICES}
\label{sec:optical_modality_families}
\label{sec:integration_architectures}

The signal paths in Section~\ref{sec:foundations_comparison} leave an important
design question: where should communication and sensing share a component,
signal, or decision? Reusing a communication carrier for sensing exposes the
measurement to its modulation and reference errors. Adding a sensing probe
provides more control over the measurement but introduces an additional signal
into the link. Moving the target interaction beyond an optical-to-radio
conversion adds another channel and calibration boundary. The appropriate
architecture therefore follows from the location of the target and the
resources through which the two functions interact.

Table~\ref{tab:modality_map} connects the platform families to these decisions.
The optical wireless families include visible light communication (VLC), light
fidelity (LiFi) networking, and free-space optical (FSO) links. Each study has
one dominant platform assignment; a system can share several
resources and contain more than one propagation medium. The discussion follows
the physical path, then examines how its shared elements become design
variables.

\begin{table*}[t]
\caption{Physical paths and architecture choices across the 206 studies}
\label{tab:modality_map}
\centering
\footnotesize
\setlength{\tabcolsep}{4pt}
\renewcommand{\arraystretch}{1.10}
\begin{tabularx}{\textwidth}{@{}>{\raggedright\arraybackslash}p{0.17\textwidth}
>{\raggedright\arraybackslash}p{0.26\textwidth}
>{\raggedright\arraybackslash}p{0.25\textwidth}
>{\raggedright\arraybackslash}X@{}}
\toprule
Platform; studies, $n$ (\%) & Target interaction and observation & Main shared
choices & Conditions that constrain the choice \\
\midrule
Photonic terahertz; 69 (33.5\%) & Optical generation followed by a wireless
millimeter wave or terahertz target path & Reference, conversion, waveform and
beam control & Source coherence, conversion bandwidth, antenna geometry and
wireless propagation \\
Fiber; 56 (27.2\%) & Guided-path perturbation, or a downstream target reached
through optical transport & Communication-carrier reuse, separate probe,
wavelength and receiver allocation & Leakage, nonlinear interaction, topology,
traffic and remote calibration \\
VLC and LiFi; 38 (18.4\%) & Direct and reflected light in an illuminated scene
& Modulation, exposure, illumination and spatial allocation & Source response,
ambient light, field of view, orientation and receiver type \\
Free space optical; 31 (15.0\%) & Directional optical link and, where present,
target return & Waveform, aperture, beam and acquisition time & Pointing,
divergence, collected power, reflectivity, weather and motion \\
Hybrid optical; 9 (4.4\%) & Interaction across optical, guided and wireless
segments & Shared reference, conversion or sensing-guided control & Segment
power, synchronization, calibration and control timing \\
Other optical; 3 (1.5\%) & Microcomb-generated microwave field, ultraviolet
response or quantum phase & Specialized source, reference and task design &
The native propagation, receiver and measurement model \\
\midrule
\textbf{Total; 206 (100.0\%)} & \multicolumn{3}{l}{One dominant family per study;
physical segments and sharing mechanisms may overlap.} \\
\bottomrule
\end{tabularx}
\end{table*}

\subsection{Guided Paths: Sensing the Fiber or Transporting the Probe}

When fiber is the sensing medium, environmental disturbances alter the same
guided field that carries traffic. Backscatter, carrier phase, modal coupling,
or an inserted sensor can reveal distributed events, vibration, strain, and
environmental change
\cite{OISAC_SCR00957,OISAC_SCR00432,OISAC_SCR00356,OISAC_SCR00679}.
The architecture first decides whether to recover this information from the
communication signal or to introduce a dedicated sensing signal. Carrier and
waveform reuse make symbols, phase, or supported modes part of the measurement
\cite{OISAC_SCR00647,OISAC_SCR00038,OISAC_SCR00015}. Separate probes or
wavelengths permit a different interrogation structure while retaining common
fiber and receiver constraints
\cite{OISAC_SCR00013,OISAC_SCR00344,OISAC_SCR00218}.

These alternatives place interference at different points. A separate probe
can leak into the data channel or induce nonlinear interaction, so its launch
conditions affect both sensing support and communication margin. Reusing the
data carrier avoids that particular addition but makes sensing depend on the
available reference and on how data modulation is removed. Coherent vibration
sensing, for example, relies on preserving and recovering phase along the
communication path \cite{OISAC_SCR00647}. Dynamic range, reflections, topology,
and traffic determine which reuse is practical; their importance persists in
live access networks and deployed routes \cite{OISAC_SCR00218,OISAC_SCR00957}.

Fiber used to transport a probe has a different measurement boundary. The
optical link distributes signals to a radio frequency identification (RFID),
radio, or radar segment, where the
target interaction occurs
\cite{OISAC_SCR00496,OISAC_SCR00631,OISAC_SCR00719,OISAC_SCR00983,OISAC_SCR01049}.
The recovered quantity may be tag state, range, localization, or an image.
Optical transport then contributes conversion, timing, and reference errors to
an observation also shaped by wireless geometry. A reported transport span
therefore cannot substitute for target range. Evaluating this architecture
requires following the outgoing signal and the return through their conversion
and clocking stages.

\subsection{Illuminated Scenes and Directional Optical Links}

In VLC and LiFi systems, the emitted light simultaneously defines the
communication path and the scene available for observation. Photodiodes favor
rapid temporal reception, while cameras provide spatially resolved samples
subject to exposure and readout constraints. Direct and reflected paths can
support position, pose, gesture, recognition, and physiological sensing
\cite{OISAC_SCR00910,OISAC_SCR01034,OISAC_SCR01036}. The receiver choice thus
changes what can be inferred from the common optical field, as well as which
modulation and update rates are usable.

Illumination makes waveform and spatial design interdependent. Nonnegativity,
dimming, flicker, and source linearity restrict temporal modulation, while
ambient light, orientation, and field of view determine which signal reaches
the detector \cite{OISAC_SCR00273,OISAC_SCR00925}. A design combining temporal
data with structured illumination uses position estimates to adapt beamforming;
the resulting coupling runs from the observed geometry back to the transmitted
field \cite{OISAC_SCR00196}. Other designs allocate emitters or spatial
resources to communication and localization
\cite{OISAC_SCR00901,OISAC_SCR00407,OISAC_SCR00557}. Here, a stronger received
signal and a more informative localization geometry are distinct objectives.
Their relation depends on emitter placement and receiver orientation.

FSO links place similar spatial decisions in a directional beam and aperture
path. The link may deliver data directly, illuminate a target, or collect a
retroreflected return. Aperture, divergence, pointing, reflectivity, and motion
then determine the available communication power and sensing observation
\cite{OISAC_SCR00011,OISAC_SCR01029,OISAC_SCR00730}. Scanning and steering
choose which region is observed and served
\cite{OISAC_SCR00012,OISAC_SCR00506,OISAC_SCR00941}; atmospheric estimates
can support beam shaping or link selection
\cite{OISAC_SCR01085,OISAC_SCR00456}. A narrow spatial response can aid
selection of a user or target, but its value depends on maintaining alignment.
Thus, the aperture and control strategy belong to the performance result,
together with the waveform. Laboratory ranging and sustained operation under
changing atmospheric conditions test different parts of this dependency
\cite{OISAC_SCR00730,OISAC_SCR01085}.

\subsection{Photonic Generation with Wireless Target Interaction}

Photonic millimeter wave and terahertz systems use optics to generate,
distribute, or process a carrier before a wireless target interaction. Laser
coherence and linewidth affect the shared reference; component bandwidth and
conversion fidelity determine the signal delivered to the antenna
\cite{OISAC_SCR00988,OISAC_SCR01004}. Antenna aperture, alignment, molecular
absorption, and target scattering subsequently shape the radiated link and
return \cite{OISAC_SCR00989,OISAC_SCR00992}. This separation identifies which
impairment a proposed photonic improvement can address and which remains in the
wireless path.

Architecture determines how far the common reference extends. A multicarrier
frame can allocate spectral support between data and sensing, while a
photonically controlled array places the joint decision in the transmitted
beam \cite{OISAC_SCR01054,OISAC_SCR00036}. Wavelength tuning or multibeam
reception moves control to the source or receiving front end
\cite{OISAC_SCR00185,OISAC_SCR00396}. Combining linear frequency modulation
(LFM) with communication symbols across coherent carriers makes waveform
generation a shared constraint \cite{OISAC_SCR00988}. Frequency hopping and
delay--Doppler multiplexing provide other ways of structuring that signal
\cite{OISAC_SCR00992,OISAC_SCR00973}; their benefits depend on how the receiver
recovers the communication symbols and target response.

The return path can introduce a further optical stage. An echo may be
remodulated onto light for centralized processing, with photonic dechirping
mapping delay to a lower-frequency observable
\cite{OISAC_SCR00966,OISAC_SCR01000}. Reference phase noise degrades recovery,
and optical/electrical path mismatch can shift the inferred range
\cite{OISAC_SCR00966}. A shared source consequently reduces neither all timing
errors nor all calibration requirements. Controlled experiments expose these
conversion and beamforming effects
\cite{OISAC_SCR00988,OISAC_SCR00036}; the reported 270~m outdoor link also
tests operation across the complete wireless chain \cite{OISAC_SCR01054}.

\subsection{Integration Across Segments and Control Loops}

Hybrid integration becomes consequential when conversion or feedback between
segments changes the system response. Fiber and wireless interactions can
couple power, timing, and calibration across a common waveform
\cite{OISAC_SCR00210,OISAC_SCR00366}. A shared photonic source can span fiber,
FSO, and millimeter wave stages \cite{OISAC_SCR00944}. In each case, an
end-to-end result depends on the interfaces as well as on the individual
segments. Separate segment measurements establish only the behavior actually
measured at those boundaries.

A control loop creates another form of integration. Sensing in one segment may
guide service in another, as in blockage-aware access and combined optical
delivery and monitoring \cite{OISAC_SCR01011,OISAC_SCR00277}. Sensed FSO
channel state can inform trajectory and resource allocation, or access
placement and user association across optical backhaul and radio access
\cite{OISAC_SCR00917,OISAC_SCR00922}. The shared element is the decision:
observation age, processing, and actuation now matter alongside signal quality.
Section~\ref{sec:technologies_applications_6g} examines the service requirements placed on
these loops.

The remaining specialized paths retain their native measurement models. A
microcomb-generated microwave field, a solar-blind ultraviolet response, and a
quantum phase measurement provide different observations and receiver
requirements \cite{OISAC_SCR00268,OISAC_SCR00379,OISAC_SCR00906}. Likewise,
integration of communication and physiological sensing in a common material
platform can retain separate detector roles \cite{OISAC_SCR00001}. Such
examples extend the architecture space without establishing equivalent
communication or sensing tasks.

\subsection{From a Shared Component to a Joint Design Choice}

Across these paths, the decisive step is to identify an adjustable choice.
Hardware sharing makes a component's bandwidth, coherence, or dynamic range
common to both functions. Carrier sharing supplies a frequency or phase
relation. Waveform sharing goes further by making a pilot, chirp, code, or
spatial pattern contribute to decoding and sensing. In OFDM, the chosen pilots
and payload structure determine which samples are known to the receiver and
how much signal support remains for data \cite{OISAC_SCR00007}. In guided
sensing, removing data modulation to recover phase makes receiver processing
part of the integration \cite{OISAC_SCR00647}.

The resulting evidence depends on how the shared choice was exercised. One
photonic prototype uses frequency steps for the radar dimension and amplitude
for data, with its experimental support strongest for waveform generation and
communication response \cite{OISAC_SCR00160}. An exposure-time design provides
an analytical and numerical coexistence result, while laboratory communication
and sensing gains were obtained in separate configurations
\cite{OISAC_SCR00273}. Another photonic system uses a shared ASK/FMCW waveform,
with functional metrics evaluated in separate receiver configurations
\cite{OISAC_SCR00083}. Localization
and imaging performed before data delivery on a shared source constitute a
further temporal arrangement \cite{OISAC_SCR00012}. These distinctions
determine which paired outcomes can be assigned to one operating point. The
next section follows the shared choice through its mechanism to those
outcomes.
